\section{Amenazas a la validez}
\label{sec:validez}
\subsection{Máquina, fechas y carga de fondo}
Los tres equipos ejecutaron procesos en localhost: comparten CPU,
memoria y loopback. Se registraron sistema, procesador, RAM y versiones,
pero no una traza continua de CPU, temperatura o programas abiertos.
Melanir y Mourad ejecutaron primero sockets y después gRPC; la serie
TCP de Steven es del 21/11, su gRPC secuencial del 19/11 y el
concurrente actualizado del 22/11. No se alternaron ni
aleatorizaron de manera controlada las condiciones. Estos factores
pueden explicar parte de las diferencias y no quedan eliminados por
tener varios miles de solicitudes.

\subsection{Qué se comparó}
La solicitud serializada tiene 100 o 10\,000 bytes; el contenido útil
no es idéntico en JSON y Protocol Buffers. El RTT de sockets excluye
la construcción del JSON, mientras la llamada gRPC incluye su
serialización interna. Son implementaciones completas con límites
distintos, no una medición aislada de TCP frente a HTTP/2
\cite{blinowski2026rpc,jackson2024streaming}.

En el programa común, el rendimiento incluye creación y cierre de
clientes y conexiones y gestión de hilos, también en el gRPC concurrente
actualizado de Steven. Esta nueva sesión usa 1\,000 envíos por cliente,
frente a 100 en las otras condiciones; una carga más larga cambia la
amortización del arranque y puede cambiar la carga de CPU y la latencia.
Sus tasas no constituyen una comparación bajo cargas equivalentes.
La serie concurrente original, con otro límite de cronometraje, queda
fuera del conjunto final y se conserva para trazabilidad.

Las conexiones de gRPC se abren con \texttt{insecure\_channel}; las
de sockets son TCP sin TLS. No se evaluaron cifrado, pérdida de paquetes,
caídas de nodos ni reintentos bajo fallo. El funcionamiento normal y
los plazos del código no prueban robustez ante fallos reales.

\subsection{Repeticiones y tratamiento de datos}
Todas las series finales tienen calentamiento etiquetado y excluido.
Se mantuvieron los identificadores de corrida y los originales, con
hashes comprobados. Cada condición secuencial tiene cinco repeticiones;
la concurrente también, incluida la sesión gRPC nueva de Steven.
Se estima la dispersión entre cinco corridas; no se trata cada
solicitud como una repetición independiente del experimento.

La desviación y P95 por mensaje describen la distribución dentro de
la muestra. Las barras de error describen una desviación entre medias
de corridas, no intervalos de confianza. No se aplicaron contrastes
de significación. Los mensajes fallidos pueden quedar fuera de los
CSV de solicitudes completadas; los metadatos no proporcionan una
medición exhaustiva de errores o pérdidas.

\subsection{Alcance y trazabilidad}
No se promedian equipos para producir una cifra única de ventaja.
Las observaciones describen máquinas, fechas y programas concretos;
una red física, la nube, TLS u otras cargas requieren nuevas pruebas.
Los CSV históricos conservan su alcance original. El registro gRPC
del 21/11 cubre también los acuses, a diferencia de las sesiones
iniciales. Los hashes permiten detectar cambios en los archivos
conservados; no sustituyen las evidencias de ejecución, la bitácora
ni la demostración del equipo.
